\subsection{Testing Requirements}

Testing shall be performed to verify that Dikaios Amanah satisfies its
functional and non-functional requirements. The testing process shall
focus on the correctness of the major system workflows, role-based
access control, data validation, security mechanisms, and system
integration.

\subsubsection{Unit Testing}

Individual software components and functions shall be tested
independently to verify that they produce the expected results.

Unit testing should cover important backend services, validation
functions, authentication-related functions, and other critical
application logic.

\subsubsection{Integration Testing}

Integration testing shall verify that different system components
communicate correctly.

The following integrations shall be considered:

\begin{itemize}
    \item Frontend and backend API communication.
    \item Backend and PostgreSQL database communication.
    \item Backend and document storage.
    \item Authentication and authorization services.
    \item Backend and the trust and audit layer.
\end{itemize}

\subsubsection{System Testing}

System testing shall evaluate the complete application as an integrated
system.

Testing shall cover the major workflows of Business Owners, Investors,
and Administrators, including business verification, campaign
management, investment simulation, profit declaration, and administrative
operations.

\subsubsection{Functional Testing}

Functional tests shall verify that the implemented system behavior
matches the functional requirements.

Representative functional test cases shall include:

\begin{itemize}
    \item User registration and login.
    \item Role-based access control.
    \item Business profile creation and management.
    \item Business verification submission.
    \item Document upload and management.
    \item Campaign creation and approval.
    \item Campaign browsing.
    \item Simulated investment.
    \item Portfolio and investment history.
    \item Profit declaration.
    \item Administrative review and monitoring.
    \item Audit record generation.
\end{itemize}

\subsubsection{Security Testing}

Security testing shall verify that protected system resources cannot be
accessed by unauthorized users.

Security testing shall include:

\begin{itemize}
    \item Authentication testing.
    \item Authorization and role testing.
    \item Invalid credential testing.
    \item Protected endpoint testing.
    \item Input validation testing.
    \item Unauthorized document access testing.
    \item Session or token validation testing.
\end{itemize}

\subsubsection{Usability Testing}

Usability testing shall evaluate whether users can understand and
complete the major system workflows using the provided interface.

The primary user roles shall be considered during usability testing,
including Business Owners, Investors, and Administrators.

\subsubsection{Performance Testing}

Basic performance testing shall be performed to evaluate the system's
response under representative academic testing conditions.

Testing may include:

\begin{itemize}
    \item API response time.
    \item Page loading behavior.
    \item Database query performance.
    \item Response behavior under multiple simulated requests.
\end{itemize}

Production-scale load and stress testing are outside the current
academic project scope.

\subsubsection{Negative Testing}

The system shall be tested with invalid, incomplete, or unexpected input
to verify that errors are handled appropriately.

Examples include:

\begin{itemize}
    \item Invalid login credentials.
    \item Missing required fields.
    \item Invalid document types.
    \item Invalid investment amounts.
    \item Unauthorized operations.
    \item Invalid campaign information.
\end{itemize}

\subsubsection{User Acceptance Testing}

User Acceptance Testing (UAT) shall verify that the system satisfies
the expected workflows from the perspective of its intended users.

Representative scenarios shall be executed for the Business Owner,
Investor, and Administrator roles.

\subsubsection{Test Documentation}

Test cases shall document the following information where applicable:

\begin{itemize}
    \item Test case identifier.
    \item Related requirement.
    \item Test objective.
    \item Preconditions.
    \item Test input.
    \item Expected result.
    \item Actual result.
    \item Test status.
\end{itemize}

Test results shall be used to identify defects and verify that corrective
changes satisfy the corresponding requirements.